\documentclass[12pt, letterpaper, landscape]{article}

% ==================================================
% PACKAGES
% ==================================================
\usepackage[margin=0.8in]{geometry}
\usepackage{xcolor}
\usepackage{pagecolor}
\usepackage{enumitem}
\usepackage{listings}
\usepackage{graphicx}
\usepackage{wrapfig}
\usepackage{caption}
\usepackage{amsmath}
\usepackage{eso-pic}

\setlength{\parindent}{0pt}
\setlength{\parskip}{0pt}

% ==================================================
% PAGE COLORS
% ==================================================
\definecolor{pagebg}{HTML}{000000}
\definecolor{textwhite}{HTML}{FFFFFF}

\pagecolor{pagebg}
\color{textwhite}

% ==================================================
% CODE COLORS
% ==================================================
\definecolor{codebg}{HTML}{181818}
\definecolor{codegreen}{HTML}{98C379}
\definecolor{codegray}{HTML}{808080}
\definecolor{codepurple}{HTML}{C678DD}
\definecolor{codeblue}{HTML}{61AFEF}
\definecolor{codewhite}{HTML}{ABB2BF}
\definecolor{codemagenta}{HTML}{E06C75}

% ==================================================
% CODE LISTING STYLE
% ==================================================
\lstdefinestyle{mystyle}{
    backgroundcolor=\color{codebg},
    basicstyle=\ttfamily\footnotesize\color{codewhite},
    commentstyle=\color{codegreen},
    keywordstyle=\color{codeblue},
    stringstyle=\color{codepurple},
    numberstyle=\tiny\color{codegray},
    numbers=left,
    numbersep=8pt,
    breaklines=true,
    breakatwhitespace=false,
    keepspaces=true,
    showspaces=false,
    showstringspaces=false,
    showtabs=false,
    tabsize=4,
    frame=single,
    rulecolor=\color{codegray},
    framerule=0.3pt,
    captionpos=b,
    xleftmargin=1.5em,
    framexleftmargin=1em,
    framesep=6pt,
    columns=fullflexible,
    upquote=true
}

\lstset{style=mystyle}

% ==================================================
% BULLET STYLE
% ==================================================
\setlist[itemize]{
    label=\textcolor{textwhite}{\textbullet},
    leftmargin=2em,
    itemsep=0.5em,
    topsep=0.5em
}

% ==================================================
% REUSABLE PAGE HEADING
% ==================================================
\newcommand{\newnotepage}[1]{
    \clearpage
    \noindent{\LARGE\bfseries #1}\par
    \vspace{0.5em}
    \noindent\rule{\textwidth}{0.5pt}
    \vspace{1em}
}

% ==================================================
% IMAGE HELPER
%
% Usage:
% \notegraphic{width}{path/to/image.png}
%
% If the image is missing, a placeholder is shown
% so the template can still compile.
% ==================================================
\newcommand{\notegraphic}[2]{%
    \IfFileExists{#2}{%
        \includegraphics[width=#1]{#2}%
    }{%
        \fcolorbox{codegray}{codebg}{%
            \parbox[c][3cm][c]{%
                \dimexpr#1-2\fboxsep-2\fboxrule\relax
            }{%
                \centering\color{textwhite}\small
                Image placeholder\\[0.5em]
                \texttt{\detokenize{#2}}%
            }%
        }%
    }%
}

% ==================================================
% DOCUMENT
% ==================================================
\begin{document}


% Title page goes here
\thispagestyle{empty}
\vspace*{\fill}
\thispagestyle{empty}
\thispagestyle{empty}
\AddToShipoutPictureBG*{%
    \AtPageLowerLeft{%
        \put(0,0){%
            \makebox(0,0)[lb]{%
                \color{black}%
                \rule{\paperwidth}{\paperheight}%
            }%
        }%
        \put(\LenToUnit{0.5\paperwidth},%
             \LenToUnit{0.5\paperheight}){%
            \makebox(0,0){%
                \includegraphics[
                    width=\paperwidth,
                    height=\paperheight,
                    keepaspectratio
                ]{test.png}%
            }%
        }%
    }%
}
\vspace*{\fill}
\begin{center}
    {\Huge\bfseries My Notes Title\par}
    \vspace{1cm}
    {\Large Subtitle or Topic\par}
    \vspace{2cm}
    {\large October 2026\par}
\end{center}
\vspace*{\fill}
\clearpage



% ==================================================
% PAGE 1: BASIC NOTES AND CODE
% ==================================================
\newnotepage{Basic Notes}

\begin{itemize}
    \item A normal bullet point.

    \item Inline subscripts: $x_1$, $V_{\text{out}}$,
    and $a_{n+1}$.

    \item An equation:
    \[
        E = mc^2
    \]

    \item Code can be inserted directly into the document.
\end{itemize}

\textbf{Example C++ code:}

\begin{lstlisting}[language=C++]
#include <iostream>

int main() {
    float x = 3.14f;
    std::cout << "Hello, world!" << std::endl;
    return 0;
}
\end{lstlisting}


% ==================================================
% PAGE 2: IMAGE ON THE RIGHT
% Text and bullets flow around the image.
% ==================================================
\newnotepage{Image on the Right}

\begin{wrapfigure}{r}{0.38\textwidth}
    \centering
    \notegraphic{0.36\textwidth}{images/diagram.png}
    \caption{Example diagram}
\end{wrapfigure}

\begin{itemize}
    \item This bullet demonstrates text wrapping beside
    an image. Add enough text to see how the lines use
    the available space to the left of the figure.

    \item The image occupies the right side of the page.
    As the text passes the bottom of the image, it can
    use the full text width again.

    \item Inline equations work here too:
    $x_{n+1} = x_n + \Delta x$.

    \item Additional notes can follow the image.
\end{itemize}


% ==================================================
% PAGE 3: IMAGE ON THE LEFT
% ==================================================
\newnotepage{Image on the Left}

\begin{wrapfigure}{l}{0.38\textwidth}
    \centering
    \notegraphic{0.36\textwidth}{images/splat.png}
    \caption{Example rendering}
\end{wrapfigure}

\begin{itemize}
    \item Use the left-hand layout to place an image
    alongside a description.

    \item This is useful for diagrams, screenshots,
    rendering results, or illustrations.

    \item Text can continue beneath the image after
    reaching the bottom edge of the figure.

    \item The figure width can be changed by editing
    the width arguments above.
\end{itemize}


% ==================================================
% PAGE 4: IMAGE BESIDE CODE
%
% Listings do not reliably wrap around wrapfigure.
% Use minipage to create two columns instead.
% ==================================================
\newnotepage{Code Beside an Image}

\textbf{Code example:}

\vspace{0.5em}

\noindent
\begin{minipage}[t]{0.57\textwidth}
    \vspace{0pt}
    \begin{lstlisting}[language=C++]
void update(float dt) {
    position += velocity * dt;
    velocity += acceleration * dt;
}
    \end{lstlisting}
\end{minipage}
\hfill
\begin{minipage}[t]{0.37\textwidth}
    \vspace{0pt}
    \centering
    \notegraphic{\linewidth}{images/code-diagram.png}
    \captionof{figure}{Update process}
\end{minipage}

\vspace{1em}

\begin{itemize}
    \item The code and image occupy separate columns.

    \item Change the minipage widths to control the
    relative amount of space each element receives.
\end{itemize}


% ==================================================
% PAGE 5: REGULAR CENTERED FIGURE
% ==================================================
\newnotepage{Standalone Figure}

\begin{figure}[htbp]
    \centering
    \notegraphic{0.7\textwidth}{images/overview.png}
    \caption{Overview diagram}
\end{figure}

\begin{itemize}
    \item Use a standalone figure when the image deserves
    most of the page width.

    \item Figures may float to a different position
    when LaTeX lays out the page.
\end{itemize}


% ==================================================
% PAGE 6: MIXED NOTES, FIGURES, AND CODE
% ==================================================
\newnotepage{Mixed Technical Notes}

\textbf{Observation}

The output depends on the input parameters
$a$, $b$, and $c$. In particular,

\[
    f(x) = \sum_{i=0}^{n} a_i x^i
\]

\begin{wrapfigure}{r}{0.34\textwidth}
    \centering
    \notegraphic{0.32\textwidth}{images/result.png}
    \caption{Result}
\end{wrapfigure}

\begin{itemize}
    \item The figure shows a representative result.

    \item The algorithm evaluates the polynomial.

    \item The output can be compared against the
    expected value.
\end{itemize}

\textbf{Implementation:}

\begin{lstlisting}[language=C++]
float evaluate(float x) {
    return a * x * x + b * x + c;
}
\end{lstlisting}

\end{document}